\documentclass[interno]{gpt}

\titolo{Decsioni a seguito dell'assegnazione del capitolato}
\data{2026-11-03}
\orario{18:00}{19:30}
\luogo{Discord}
\redattore{mazzurana}
\verificatore{corro}
%\assenti{spagnoli}

%\modifica{1.0.0}{2026-10-29}{corro}{-}{Verifica e approvazione}
\modifica{0.1.0}{2026-11-03}{mazzurana}{-}{Prima stesura}

\componenti{biasini,cescobolla,corro,guglielmo,mazzurana,spagnoli}
\begin{document}

\section{Ordine del giorno}
\begin{enumerate}
  \item Assegnazione ruoli di progetto per il prossimo sprint.
  \item Valutazione e risoluzione delle criticità segnalate dal docente (sito documentazione e versionamento).
  \item Pianificazione della stesura di analisi dei requisiti, piano di progetto e piano di qualifica.
  \item Pianificazione incontro con il proponente Eggon.
\end{enumerate}

\clearpage

\section{Svolgimento}
\subsection{Definizione e rotazione dei ruoli di progetto}
In apertura di riunione il gruppo ha definito l'assegnazione e la pianificazione dei ruoli per il periodo iniziale:
\begin{itemize}
  \item \textbf{Prima settimana}: considerata la priorità assoluta per l'avvio dell'analisi dei requisiti e la temporanea non operatività di programmatore e progettista, i ruoli operativi sono stati provvisoriamente distribuiti aumentando la copertura sull'analisi:
  \begin{itemize}
    \item Amministratore: Enrico Biasini
    \item Responsabile: Luca Cesco Bolla
    \item Analista: Gaetano Spagnoli
    \item Analista: Francesco Guglielmo
    \item Analista: Francesco Mazzurana
    \item Verificatore: Alessio Corrò
  \end{itemize}
  \item \textbf{Seconda settimana}: ripristino dell'assetto standard a copertura di tutti i ruoli formali di progetto:
  \begin{itemize}
    \item Amministratore: Enrico Biasini
    \item Responsabile: Luca Cesco Bolla
    \item Analista: Gaetano Spagnoli
    \item Progettista: Francesco Guglielmo
    \item Programmatore: Francesco Mazzurana
    \item Verificatore: Alessio Corrò
  \end{itemize}
\end{itemize}
Si è inoltre concordato che, a partire dalla configurazione standard, la rotazione avverrà a scorrimento in avanti secondo un'indicizzazione modulo 6 (chi ricopre il ruolo di Amministratore diverrà Responsabile, il Responsabile diverrà Analista, e così via). Tale schema funge da modello guida ed è suscettibile di eventuali adattamenti in caso di imprevisti o esigenze di carico di lavoro.
\subsection{ Valutazioni e decisioni a seguito dei feedback ricevuti dal docente}
A seguito dei feedback ricevuti dal docente nell'aggiudicazione dell'appalto, si provvede ad intervenire per risolvere le criticità emerse lato vista web per la documentazione e problematiche del versionamento.
\subsection{Ripartizione dei compiti per la documentazione}
Per procedere in modo efficace verso i prossimi traguardi, sono state suddivise le responsabilità per la redazione della documentazione di ingresso: 
\begin{itemize}
  \item \textbf{analisi dei requisiti:} affidata a Enrico Biasini e Luca Cesco Bolla. 
  \item \textbf{piano di progetto:} affidato a Gaetano Spagnoli e Francesco Guglielmo. 
  \item \textbf{piano di qualifica:} affidato a Francesco Mazzurana e Alessio Corrò. 
\end{itemize}
Ciascuna coppia curerà l'impostazione iniziale della propria sezione coordinandosi con il verificatore e l'amministratore per garantire coerenza.

\subsection{Contatto con il proponente Eggon}
È stata inviata una comunicazione formale per concordare data, orario e piattaforma telematica per il primo meeting di allineamento sui requisiti e sulle aspettative del capitolato d'appalto con la proponente.


\section{Decisioni}
\begin{decisioni}
  \decisione{Assegnazione formale dei ruoli per il ciclo corrente.}
  \decisione{Risoluzione prioritaria delle anomalie segnalate dal docente relative al portale della documentazione e al versionamento.}
  \decisione{Suddivisione a coppie per l'inizio della stesura dei documenti analisi dei requisiti, piano di progetto e piano di qualifica.}
\end{decisioni}

\section{Azioni da intraprendere}
\begin{azioni}
  \azione[\#65]{Stesura verbale interno della riunione odierna.}{mazzurana}{03/11/2026}
  \azione[\#67]{Adeguamento versionamento e portale della documentazione.}{spagnoli, guglielmo}{04/11/2026}
  \azione[\#68]{Approfondimento e inizio stesura dei documenti analisi dei requisiti, piano di progetto e piano di qualifica}{tutti}{05/11/2026}
  \azione[\#71]{Contatto con Eggon per organizzazione incontro}{cescobolla}{03/11/2026}
\end{azioni}

\section{Prossima riunione}
La prossima riunione interna è fissata per \textbf{giovedì 5 novembre 2026 alle ore 18:00} sulla piattaforma Discord. L'obiettivo principale dell'incontro sarà legato ai documenti da produrre, ovvero: analisi dei requisiti, piano di progetto e piano di qualifica.

\end{document}
